\SourceLeaf{259}%
\(DB\). Für diesen Punkt bleibt \(p\) gleich, und die Zeichen \(q',r',\varphi',\psi',S'\) bezeichnen dasselbe wie \(q,r,\varphi,\psi,S\) für den Punkt \(B\). So entsteht ein Dreieck mit den Eckpunkten \(A,B,C\). Seine Winkel bezeichnen wir mit \(A,B,C\), die gegenüberliegenden Seiten mit \(a,b,c\) und seinen Flächeninhalt mit \(\sigma\). Das Krümmungsmaß in den Punkten \(A,B,C\) bezeichnen wir der Reihe nach mit \(\alpha,\beta,\gamma\). Nehmen wir also, wie es zulässig ist, an, dass \(p,q,q-q'\) positiv sind, so gilt
\[
\begin{gathered}
A=\varphi-\varphi',\quad B=\psi,\quad C=\pi-\psi',\\
a=q-q',\quad b=r',\quad c=r,\quad\sigma=S-S'
\end{gathered}
\]

Zunächst drücken wir den Flächeninhalt \(\sigma\) durch eine Reihe aus. Ersetzen wir in [7] jede auf \(B\) bezogene Größe durch die entsprechende Größe für \(C\), so ergibt sich die Formel für \(S'\). Daraus erhalten wir bis zu Größen sechster Ordnung
\[
\begin{aligned}
\sigma=\tfrac{1}{2}p(q-q')\{&1-\tfrac{1}{6}f^0(pp+qq+qq'+q'q')\\
&-\tfrac{1}{60}f'p(6pp+7qq+7qq'+7q'q')\\
&-\tfrac{1}{20}g^0(q+q')(3pp+4qq+4qq'+4q'q')\}
\end{aligned}
\]
Mithilfe der Reihe [2], nämlich
\[
c\sin B=p(1-\tfrac{1}{3}f^0qq-\tfrac{1}{4}f'pqq-\tfrac{1}{2}g^0q^3-\text{etc.})
\]
geht diese Formel in die folgende über
\[
\begin{aligned}
\sigma=\tfrac{1}{2}ac\sin B\{&1-\tfrac{1}{6}f^0(pp-qq+qq'+q'q')\\
&-\tfrac{1}{60}f'p(6pp-8qq+7qq'+7q'q')\\
&-\tfrac{1}{20}g^0(3ppq+3ppq'-6q^3+4qqq'+4qq'q'+4q'^3)\}
\end{aligned}
\]

Das Krümmungsmaß in einem beliebigen Punkt der Fläche ergibt sich (nach Artikel 19, in dem \(m,p,q\) das bezeichneten, was wir hier mit \(n,q,p\) bezeichnen)
\[
\begin{aligned}
&=-\frac{1}{n}\mathord{.}\frac{\mathrm{d}\mathrm{d}n}{\mathrm{d}q^2}=-\frac{2f+6gq+12hqq+\text{etc.}}{1+fqq+\text{etc.}}\\
&=-2f-6gq-(12h-2ff)qq-\text{etc.}
\end{aligned}
\]
Daraus folgt, wenn sich \(p,q\) auf den Punkt \(B\) beziehen,
\[
\beta=-2f^0-2f'p-6g^0q-2f''pp-6g'pq-(12h^0-2f^0f^0)qq-\text{etc.}
\]
sowie
\[
\begin{aligned}
\gamma&=-2f^0-2f'p-6g^0q'-2f''pp-6g'pq'-(12h^0-2f^0f^0)q'q'-\text{etc.}\\
\alpha&=-2f^0
\end{aligned}
\]
